\documentclass[10pt]{article}
\usepackage[utf8]{inputenc}
\usepackage{amsmath}
\usepackage{geometry}
\usepackage{enumitem}
\geometry{a4paper, margin=0.5in}

\setlist[enumerate]{itemsep=1pt, topsep=2pt, parsep=0pt}

\begin{document}

\begin{center}
    \textbf{\large EXAMEN PARCIAL DE FÍSICA BÁSICA (FIS-0001) -- TEMA B} \\
    \textbf{INSTITUTO TECNOLÓGICO SUPERIOR COMUNITARIO (ITSC)} \\
    \textbf{Docente: Lic. Carlos Ramírez \hfill Calificación: \rule{1.5cm}{0.1pt} / 10 Puntos}
\end{center}

\vspace{0.1cm}
\noindent Nombre: \hrulefill \ Matrícula: \hrulefill \ Sección: \hrulefill \ Fecha: \hrulefill

\vspace{0.2cm}
\begin{center}
    \framebox[\textwidth]{\parbox{0.95\textwidth}{\centering \small Fórmulas de apoyo:
        $t = \frac{d}{v} \quad A_y = A \operatorname{sen}\theta \quad d = v \cdot t \quad A = \sqrt{A_x^2 + A_y^2} \quad \vec{R} = (A_x + B_x, A_y + B_y) \quad \tan\theta = \frac{A_y}{A_x} \quad A_x = A \cos\theta \quad v = \frac{d}{t}$
    }}
\end{center}

\vspace{0.2cm}
\noindent \textbf{PARTE I: CUESTIONARIO TEÓRICO (4 Puntos -- 0.5 pto c/u)} \\
\textit{Instrucciones: Encierre en un círculo la letra que represente la opción correcta para cada enunciado.}

\begin{enumerate}[leftmargin=*]
    \item ¿Cuál de los siguientes enunciados justifica de forma correcta el uso de la notación científica? \\
    a) Permite escribir de manera simplificada cantidades sumamente grandes o muy pequeñas mediante potencias de 10. \\
    b) Garantiza que no se cometan errores al realizar sumas aritméticas decimales. \\
    c) Redondea automáticamente todas las cifras significativas a números enteros.

    \item ¿Cuántas cifras significativas tiene la cantidad medida de voltaje $0.009040\text{ V}$? \\
    a) Tres cifras significativas \hfill b) Cuatro cifras significativas \hfill c) Seis cifras significativas \hfill \mbox{}

    \item Al restar de la medición $24.6\text{ m}$ la longitud $3.18\text{ m}$, el resultado expresado con el número correcto de cifras significativas es: \\
    a) $21.42\text{ m}$ \hfill b) $21.4\text{ m}$ \hfill c) $21\text{ m}$ \hfill \mbox{}

    \item ¿Cuál de las siguientes cantidades físicas representa una magnitud escalar? \\
    a) El desplazamiento \hfill b) El peso \hfill c) La energía \hfill \mbox{}

    \item Físicamente, la dirección de un vector en el plano cartesiano está determinada por: \\
    a) La longitud escalar del segmento de recta dirigida. \\
    b) El ángulo de inclinación de la línea de acción respecto al eje horizontal positivo. \\
    c) El sentido indicado por la punta de la flecha del vector.

    \item Si dos vectores de velocidad en el plano son ortogonales (perpendiculares) entre sí, el módulo de su vector resultante se calcula mediante: \\
    a) La diferencia aritmética de los módulos individuales. \\
    b) El Teorema de Pitágoras sobre las magnitudes vectoriales. \\
    c) El cociente directo entre ambos valores escalares.

    \item La longitud total medida sobre el camino real que describe un cuerpo en movimiento se define como: \\
    a) Distancia recorrida \hfill b) Desplazamiento neto \hfill c) Posición instantánea \hfill \mbox{}

    \item En una gráfica de posición contra tiempo ($x$ vs $t$) para un objeto en Movimiento Rectilíneo Uniforme (M.R.U.), la velocidad está representada por: \\
    a) La curvatura o área bajo la curva del gráfico. \\
    b) La pendiente de la línea recta resultante. \\
    c) La altura de la línea horizontal con respecto al eje de tiempo.
\end{enumerate}

\vspace{0.3cm}
\noindent \textbf{PARTE II: EJERCICIOS PRÁCTICOS (6 Puntos -- 1.5 ptos c/u)} \\
\textit{Instrucciones: Resuelva de forma ordenada cada ejercicio. Los problemas deben ser justificados en la parte de atrás de esta hoja. Escriba el resultado final en la línea correspondiente.}

\begin{enumerate}[leftmargin=*]
    \item Realice la siguiente operación y exprese el resultado final en notación científica: \\
    $(2.00 \times 10^6\text{ m/s}) \times (3.00 \times 10^{-3}\text{ s})$ \\
    \textbf{Respuesta:} Resultado = \rule{5cm}{0.1pt} m

    \item Efectúe la siguiente suma de mediciones y exprese el resultado final redondeado al número correcto de cifras significativas: \\
    $14.28\text{ s} + 3.3\text{ s} + 0.245\text{ s}$ \\
    \textbf{Respuesta:} Suma = \rule{5cm}{0.1pt} s

    \item Un vector velocidad $\vec{v}$ tiene una magnitud de $10.00\text{ m/s}$ y forma un ángulo de $45.00^\circ$ con el eje horizontal positivo. Determine los valores de sus componentes rectangulares $v_x$ y $v_y$. (Dato: $\cos 45.00^\circ = \operatorname{sen} 45.00^\circ \approx 0.71$) \\
    \textbf{Respuesta:} $v_x = $ \rule{3cm}{0.1pt} m/s \quad | \quad $v_y = $ \rule{3cm}{0.1pt} m/s

    \item Una motocicleta viaja por una avenida recta manteniendo una velocidad constante de $15.00\text{ m/s}$. Calcule la distancia en metros que logrará recorrer al cabo de un intervalo de tiempo de $4.00\text{ minutos}$ ($240.00\text{ segundos}$). \\
    \textbf{Respuesta:} Distancia = \rule{5cm}{0.1pt} m
\end{enumerate}

% CLAVE DE RESPUESTAS PARA CONTROL DEL DOCENTE (SÓLO CÓDIGO FUENTE):
% PARTE I (Teoría):
% 1. a
% 2. b (Cifras significativas: 9, 0, 4, 0)
% 3. b (Se redondea a un decimal debido a 24.6: 21.42 -> 21.4)
% 4. c (La energía es escalar; peso y desplazamiento son vectoriales)
% 5. b
% 6. b
% 7. a
% 8. b
% 
% PARTE II (Ejercicios):
% 1. 6.00 \times 10^3 m (o 6000.00 m)
% 2. 17.8 s (Suma matemática da 17.825, pero se limita a un decimal por el sumando "3.3")
% 3. v_x \approx 7.10 m/s (10.00 * 0.71) | v_y \approx 7.10 m/s (10.00 * 0.71)
% 4. 3600.00 m o 3.60 \times 10^3 m (d = v * t = 15.00 * 240.00)

\end{document}
